% Overleaf 编译器请选择 XeLaTeX。
% fontset=fandol 使用 TeX Live 自带中文字体，不依赖本机宋体、黑体或苹方。
\documentclass[UTF8,fontset=fandol,10pt,a4paper]{ctexart}
\usepackage[
  top=1.55cm,
  bottom=1.55cm,
  left=1.70cm,
  right=1.70cm
]{geometry}
\usepackage[dvipsnames]{xcolor}
\usepackage{enumitem}
\usepackage{tabularx}
\usepackage{array}
\usepackage{titlesec}
\usepackage{hyperref}
\usepackage{microtype}
\usepackage{lastpage}
\usepackage{fancyhdr}

\definecolor{primary}{RGB}{0,79,144}
\definecolor{secondary}{RGB}{72,72,72}
\definecolor{rulecolor}{RGB}{185,195,205}

\hypersetup{
  colorlinks=true,
  urlcolor=primary,
  pdftitle={陆书鸿 - AI应用与Agent工程师简历},
  pdfauthor={陆书鸿},
  pdfcreator={XeLaTeX}
}

\pagestyle{fancy}
\fancyhf{}
\renewcommand{\headrulewidth}{0pt}
\renewcommand{\footrulewidth}{0pt}
\fancyfoot[C]{\color{secondary}\footnotesize 陆书鸿 \quad · \quad 第 \thepage\ 页 / 共 \pageref*{LastPage} 页}

\setlength{\parindent}{0pt}
\setlength{\parskip}{0pt}
\setlength{\tabcolsep}{0pt}
\setlist[itemize]{
  leftmargin=1.35em,
  label=\textcolor{primary}{\small\textbullet},
  itemsep=0.30em,
  topsep=0.32em,
  parsep=0pt,
  partopsep=0pt
}

\titleformat{\section}
  {\large\bfseries\color{primary}}
  {}{0pt}{}
  [\vspace{-0.25em}\color{rulecolor}\titlerule]
\titlespacing*{\section}{0pt}{0.85em}{0.45em}

\newcommand{\daterange}[1]{\textcolor{secondary}{\small #1}}
\newcommand{\companyentry}[4]{%
  \begin{tabularx}{\textwidth}{@{}X>{\raggedleft\arraybackslash}p{3.1cm}@{}}
    \textbf{#1} \quad \textbf{#2} & \daterange{#3} \\
    \multicolumn{2}{@{}l@{}}{\textcolor{secondary}{\small #4}}
  \end{tabularx}
}
\newcommand{\projecttitle}[2]{%
  \vspace{0.18em}%
  \begin{tabularx}{\textwidth}{@{}X>{\raggedleft\arraybackslash}p{7.2cm}@{}}
    \textbf{#1} & \textcolor{secondary}{\small #2}
  \end{tabularx}
}
\newcommand{\educationentry}[3]{%
  \begin{tabularx}{\textwidth}{@{}X>{\raggedleft\arraybackslash}p{3.1cm}@{}}
    \textbf{#1} \quad #2 & \daterange{#3}
  \end{tabularx}
}

\begin{document}

\begin{center}
  {\fontsize{22pt}{24pt}\selectfont\bfseries 陆书鸿}\\[0.18em]
  {\large\bfseries\color{primary} AI应用研发工程师｜Agent \& Applied ML}\\[0.32em]
  \small
  杭州
  \quad\textcolor{rulecolor}{|}\quad
  15868855774
  \quad\textcolor{rulecolor}{|}\quad
  \href{mailto:415400396@qq.com}{415400396@qq.com}
  \quad\textcolor{rulecolor}{|}\quad
  \href{https://github.com/415400396}{github.com/415400396}
\end{center}

\vspace{-0.35em}

\section{个人简介}

具备3年以上AI/ML研发经历，兼具生产落地经验，南加州大学计算机科学硕士。经历覆盖生成模型研究、模型框架、业务机器学习与LLM应用，近一年负责换电时效预测和车辆阈值智能调整，从数据与模型评测延伸到Java服务、结构化解析、领域规则、人工审批和审计。擅长把不确定的模型输出转化为可验证、可约束的业务系统；IROS 2023第一作者、BMVC 2023第二作者。

\section{核心能力}

\begin{tabularx}{\textwidth}{@{}>{\bfseries}p{2.45cm}X@{}}
LLM与Agent应用 & Prompt Engineering、Structured Output、领域实体增强、LLM离线评测、确定性校验、人工审批（HITL）与受控Agent接入 \\
Applied ML与数据 & XGBoost、PyTorch、特征工程、日期OOT、同样本/分层评测、Spark/Hive、StarRocks、线上线下一致性验证 \\
工程落地 & Python、Java、SQL、Spring Boot、Streamlit、REST API、MySQL、Linux、Git
\end{tabularx}

\section{职业经历}

\companyentry
  {滴滴出行（杭州青桔科技有限公司）}
  {高级算法工程师（D6）}
  {2025.09 -- 至今}
  {两轮车业务线 · 电单车换电方向 · 用户保障策略组}

\projecttitle{全国换电工单时效预测}{XGBoost · Spark/Hive · Java · StarRocks}
\begin{itemize}
  \item 接手既有全国换电工单时效系统，负责后续模型与工程演进；审计181日、近783万行时序明细，重新核验混合事件Label，构建日期OOT与时点一致特征，收敛为87维XGBoost候选模型。
  \item 在318,499条、228城公平同工单集上，将MAE从126.84分钟降至123.91分钟（-2.30\%），Spearman由0.4965提升至0.5280，196/228城改善；通过分桶评测识别收益主要来自2小时以上样本，并保留短时段退化的业务边界。
  \item 建立同工单、分桶、分布标准化及特征消融评测，识别事件语义迁移造成的虚假提升；在XGBoost 0.90 Java约束下冻结87维、536树CPU模型与特征契约，99/99项检查通过，同一向量线上/本地预测差异约0.399分钟。
\end{itemize}

\projecttitle{车辆阈值智能调整与Agent接入}{LLM · Java/Spring Boot · BPM · MySQL}
\begin{itemize}
  \item 负责车辆阈值智能调整核心链路，将运营自然语言转为受Schema约束的领域命令；以权威实体知识限定城市/运营区范围，由数据库提供当前配置，并用Java校验字段、取值、动作互斥与批量完整性，人工路径经影响预览、确认、BPM、确定性执行和审计完成高风险修改。
  \item 建设长期基础、周期调整和日期特例的分层配置语义，支持相对/绝对调整、复制和固定阈值等操作；采用完整读改写避免局部车型修改覆盖同一配置中的其他字段，并将外部Agent结构化入口与人工入口隔离。
  \item 基于42条领域金标、239次模型调用完成离线评测：候选模型裸通过40/42，经Java后置校验为41/42、0截断；43条真实请求离线回放中平均延迟由原日志34.43秒降至7.17秒（-79.2\%），但一次业务可接受仅40/43，未因速度提升直接替换，而是增加批量覆盖检查与失败阻断。
  \item 基于既有安全执行底座，完成城市异动诊断Agent演进方案设计：以受控只读工具和状态化Workflow编排数据质量、异常检测、证据下钻、确定性建议、Dry Run、人审、执行终态和效果回看；该诊断Workflow尚未上线。
\end{itemize}

\newpage

\section{职业经历（续）}

\companyentry
  {中国电子科技南湖研究院（中电海康）}
  {算法工程师}
  {2024.05 -- 2025.08}
  {大模型应用、深度学习框架与模型部署}

\projecttitle{基于CBT流程的多智能体对话数据生成系统}{Qwen-32B · Unsloth · RAG · SFT}
\begin{itemize}
  \item 基于约2,000条CBT心理咨询流程数据与公开知识库，设计患者、咨询师、用户分析和流程控制角色，将咨询拆分为信息收集、问题识别、认知干预和总结反馈等可控阶段。
  \item 按阶段组织上下文、角色目标与Prompt，构建RAG、多角色协作、规则校验和人工抽检的数据生成流水线；使用Unsloth对Qwen-32B进行分角色、分阶段SFT，产出心理垂直领域多轮对话训练语料。
\end{itemize}

\projecttitle{大模型与深度学习框架研发}{JAX · Qwen MoE · CUDA/Triton · vLLM}
\begin{itemize}
  \item 参与Qwen MoE架构重构与推理适配验证，探索基于SNN阈值的动态专家路由，并开发CUDA/Triton自定义算子与vLLM适配原型。
  \item 基于JAX建设模型转换与部署工具链，实现JAX到ONNX的端到端转换，完成LLaMA、BERT、YOLO等模型的算子兼容和部署验证。
\end{itemize}

\vspace{0.25em}
\companyentry
  {南加州大学创意技术研究所（USC ICT）}
  {全职研究员}
  {2023.01 -- 2024.02}
  {生成模型、人体动作生成与视频理解}

\begin{itemize}
  \item 研究基于VQ/RQ-VAE与离散生成模型的语音手势和文本动作生成，以离散动作表示建模一对多映射；参与GENEA Challenge 2022并发表ICMI 2022论文，后续完成IROS 2023第一作者工作。
  \item 参与文本到动作离散扩散研究并发表BMVC 2023论文；设计用户研究与Web评估平台，联合主观评测和客观指标验证生成质量。
  \item 基于MMAction2开展群体活动识别，复现DIN与ST Transformer，融合Skeleton/RGB特征，并探索Cross-Attention与Domain Adaptation缓解合成数据域偏差。
\end{itemize}

\section{精选论文}

\begin{itemize}
  \item \href{https://doi.org/10.48550/arXiv.2303.12822}{\textbf{Co-Speech Gesture Synthesis using Discrete Gesture Token Learning}}，IROS 2023，\textbf{第一作者}。
  \item \href{https://proceedings.bmvc2023.org/624/}{\textbf{Text-to-Motion Synthesis using Discrete Diffusion Model}}，BMVC 2023，\textbf{第二作者}。
  \item \href{https://doi.org/10.1145/3610661.3616556}{\textbf{Discrete Diffusion for Co-Speech Gesture Synthesis}}，ICMI 2023。
  \item \href{https://doi.org/10.1145/3536221.3558059}{\textbf{The DeepMotion Entry to the GENEA Challenge 2022}}，ICMI 2022。
\end{itemize}

\section{教育经历}

\educationentry{南加州大学}{计算机科学硕士；GPA 3.83/4.0；Master's Student Honors Program}{2020.08 -- 2022.12}
\vspace{0.18em}
\educationentry{浙江工业大学}{计算机科学与技术学士；GPA 4.03/5.0，排名2/85}{2016.09 -- 2020.06}
\begin{itemize}
  \item 三次浙江省政府奖学金（Top 5\%）；全国大学生数学竞赛浙江赛区二等奖；2019年加州大学伯克利分校交换学习。
\end{itemize}

\section{技术技能}

\begin{tabularx}{\textwidth}{@{}>{\bfseries}p{2.4cm}X@{}}
编程与工程 & Python、Java、C++、SQL；Spring Boot、Streamlit、REST API、MySQL、Linux、Git \\
大模型与Agent & Qwen、Prompt Engineering、Structured Output、RAG、Agent Workflow设计、LLM Eval、HITL、Unsloth、vLLM \\
机器学习与数据 & PyTorch、XGBoost、JAX、Transformer、VAE、Diffusion Model、Spark/Hive、StarRocks
\end{tabularx}

\end{document}
